\documentclass[11pt]{article}
\usepackage[utf8]{inputenc}
\usepackage[T1]{fontenc}
\usepackage{amsmath,amssymb,amsthm,mathtools}
\usepackage[margin=1in]{geometry}
\usepackage{enumitem}
\usepackage{xcolor}

\newcommand{\PR}{\textcolor{green!55!black}{\textbf{[P]}}}
\newcommand{\OP}{\textcolor{red!75!black}{\textbf{[O]}}}
\newcommand{\GAP}{\textcolor{red!75!black}{\textsc{[gap]}}}

\theoremstyle{plain}
\newtheorem{theorem}{Teorema}
\newtheorem{lemma}{Lema}
\newtheorem{proposition}{Proposici\'on}
\newtheorem{corollary}{Corolario}
\theoremstyle{definition}
\newtheorem{remark}{Observaci\'on}
\newtheorem{definition}{Definici\'on}

\newcommand{\R}{\mathbb{R}}
\newcommand{\half}{\tfrac12}
\newcommand{\El}{E_\lambda}
\newcommand{\Sframe}{\mathcal S_\lambda}
\newcommand{\Aweil}{\mathcal A_\lambda}
\DeclareMathOperator{\spec}{Spec}

\title{\textbf{Objetivo B, veredicto: densidad vs.\ posici\'on.\\
El $\varepsilon_0$ incondicional es el operador de \emph{frame} (densidad);\\
RH vive en la brecha frame$\,\ne\,$Weil (posici\'on)}}
\author{Programa \texttt{riemann-program} --- documento de trabajo}
\date{\today}

\begin{document}
\maketitle

\begin{abstract}
La pregunta que decid\'ia si el programa CCM da un puente \emph{no circular} a RH era:
¿la cota inferior $\varepsilon_0(\lambda)\ge c/\lambda^2$ se sigue de la \emph{densidad}
de los ceros (PNT/Landau, no circular) o de su \emph{posici\'on} (RH)? La respuesta,
una vez separados con cuidado dos operadores que el programa ven\'ia confundiendo, es
n\'itida: \textbf{el $\varepsilon_0$ probado incondicionalmente es el menor autovalor
del operador de \emph{frame} $\Sframe$ (puro de densidad, Landau--Beurling), y RH
\emph{no} es ese objeto: RH es la coincidencia $\Aweil=\Sframe$ (la forma de Weil
iguala al frame), que es exactamente la condici\'on de \emph{posici\'on}.} El programa
entrega resultados de densidad genuinos e incondicionales; RH queda \emph{fuera} de su
alcance no circular, en la brecha $\Aweil-\Sframe$, que es RH por definici\'on (criterio
de Weil). Veredicto: reformulaci\'on limpia, no puente.
\end{abstract}

\tableofcontents

% =====================================================================
\section{Dos operadores que hay que separar}\label{sec:two}
% =====================================================================

Sea $\El=\mathrm{PW}_{L/2}$, $L=2\log\lambda$, y $\{\gamma_j\}$ las alturas de los
ceros no triviales $\rho_j=\beta_j+i\gamma_j$ de $\zeta$ con $\gamma_j$ bajo el
horizonte $T^*\sim2\pi\lambda^2$.

\begin{definition}[Operador de frame $\Sframe$ --- densidad]
\[
  \Sframe(f,f)\;:=\;\sum_{j}\bigl|\widehat f(\half+i\gamma_j)\bigr|^2,
  \qquad
  \varepsilon_0^{\Sframe}(\lambda):=\min_{f\in\El,\;\|f\|=1}\Sframe(f,f).
\]
Es la forma de muestreo de $\widehat f$ en las \emph{alturas reales} $\gamma_j$.
Es $\ge0$ por construcci\'on (suma de cuadrados); $>0$ sii $\{\gamma_j\}$ es un
\emph{frame} para $\El$, lo cual depende \emph{s\'olo de la densidad} de las
$\gamma_j$ (teorema de densidad de Landau/Beurling).
\end{definition}

\begin{definition}[Forma de Weil $\Aweil$ --- posici\'on]
\[
  \Aweil(f,f)\;:=\;\sum_{j}\widehat f(\rho_j)\,\overline{\widehat f(1-\overline{\rho_j})}
  \;=\;\sum_j \widehat f(\beta_j+i\gamma_j)\,\overline{\widehat f(1-\beta_j+i\gamma_j)} .
\]
Es la forma que aparece en la f\'ormula expl\'icita; involucra las \emph{posiciones
completas} $\rho_j$ (parte real $\beta_j$ incluida).
\end{definition}

% =====================================================================
\section{La identidad clave: $\Aweil=\Sframe \Leftrightarrow$ RH}\label{sec:key}
% =====================================================================

\begin{lemma}[Coincidencia frame--Weil]\label{lem:coincide}
Para un cero individual $\rho_j$:
\[
  \widehat f(\rho_j)\,\overline{\widehat f(1-\overline{\rho_j})}
  \;=\;\bigl|\widehat f(\half+i\gamma_j)\bigr|^2
  \quad\Longleftrightarrow\quad
  \beta_j=\half .
\]
En efecto, $\rho_j=\half+i\gamma_j$ da $1-\overline{\rho_j}=\half+i\gamma_j=\rho_j$,
y el t\'ermino se vuelve $|\widehat f(\rho_j)|^2$. Rec\'iprocamente, si
$\beta_j\ne\half$ entonces $\rho_j\ne1-\overline{\rho_j}$ y el t\'ermino es un producto
cruzado $\widehat f(\rho_j)\overline{\widehat f(1-\overline{\rho_j})}$, en general
\emph{no} real $\ge0$.
\end{lemma}

\begin{theorem}[El $\varepsilon_0$ del programa es densidad; RH es posici\'on]
\label{thm:verdict}
\begin{enumerate}[label=\textbf{(\alph*)},leftmargin=2em]
\item \PR\ \emph{(Densidad, incondicional.)} $\varepsilon_0^{\Sframe}(\lambda)>0$ para
todo $\lambda$, con $\varepsilon_0^{\Sframe}(\lambda)\sim C_0/\lambda^2$ (d\'eficit de
muestreo cr\'itico, Landau--Beurling). \emph{Esto es lo que el programa midi\'o y
``prob\'o''}; no usa la posici\'on de los ceros, s\'olo su densidad
$\tfrac1{2\pi}\log\tfrac{T}{2\pi}$.
\item $\Aweil(f,f)=\Sframe(f,f)$ para todo $f\in\El$ y todo $\lambda$
$\;\Longleftrightarrow\;$ $\beta_j=\half$ para todo $\gamma_j<T^*(\lambda)$
$\;\Longleftrightarrow\;$ RH bajo el horizonte.
\item \OP\ \emph{(Posici\'on $=$ RH.)} $\Aweil(f,f)\ge0$ para todo $f\in\El$
$\;\Longleftrightarrow\;$ RH bajo el horizonte. La cota inferior
$\varepsilon_0^{\Aweil}(\lambda)\ge c/\lambda^2$ es por tanto RH--equivalente.
\end{enumerate}
\end{theorem}
\begin{proof}
(a) es el teorema de densidad de Landau aplicado al frame $\{\gamma_j\}$ a densidad
cr\'itica $L/2\pi$; la tasa $C_0/\lambda^2$ es el d\'eficit en el borde
$T^*\sim2\pi\lambda^2$. (b) es el Lema \ref{lem:coincide} sumado sobre $j$ (los
t\'erminos coinciden uno a uno sii cada $\beta_j=\half$). (c): si todos $\beta_j=\half$
entonces $\Aweil=\Sframe\ge0$; si alg\'un $\beta_{j_0}\ne\half$ con $\gamma_{j_0}<T^*$,
se elige $f$ concentrada espectralmente cerca de $\rho_{j_0}$ y $1-\overline{\rho_{j_0}}$
con fases opuestas, haciendo $\Aweil(f,f)<0$; por homogeneidad
$\varepsilon_0^{\Aweil}\le0$. \qed
\end{proof}

% =====================================================================
\section{Por qu\'e la confusi\'on import\'o}\label{sec:confusion}
% =====================================================================

\begin{remark}[La fuente del espejismo]
A lo largo del programa, $\varepsilon_0$ se trat\'o como ``el margen de positividad de
Weil'' \emph{y} como ``el menor autovalor del frame de los ceros'', como si fueran el
mismo objeto. El Teorema \ref{thm:verdict} muestra que coinciden \emph{exactamente bajo
RH} (Lema \ref{lem:coincide}) y \emph{s\'olo} entonces. Como toda la num\'erica usa los
ceros \emph{verdaderos} (que est\'an en l\'inea hasta alturas enormes), se midi\'o
siempre el r\'egimen $\Aweil=\Sframe$, y el $\varepsilon_0\sim C_0/\lambda^2$ observado
es el del \emph{frame} (densidad). El contenido RH --- la posibilidad
$\Aweil\ne\Sframe$ --- es \emph{invisible} a esa medici\'on, porque requerir\'ia un cero
fuera de l\'inea, que no existe en el rango computado.
\end{remark}

\begin{corollary}[Qu\'e da la densidad y qu\'e no]
\begin{itemize}[leftmargin=1.4em]
\item \PR\ La densidad (PNT/Landau) da $\varepsilon_0^{\Sframe}\sim C_0/\lambda^2>0$:
real, incondicional, pero es una propiedad de \emph{muestreo}, no de RH.
\item Las regiones libres de ceros cl\'asicas (Vinogradov--Korobov) dan una cota
\emph{parcial} y \emph{no uniforme} sobre $\beta_j$ ($\beta_j<1-c/(\log\gamma_j)^{2/3}$),
que controla $\Aweil-\Sframe$ \emph{s\'olo} en una franja que se angosta; \textbf{no}
alcanza $\beta_j=\half$.
\item \OP\ El salto de ``franja que se angosta'' a $\beta_j=\half$ uniforme es
exactamente RH. Ninguna entrada de densidad lo cubre.
\end{itemize}
\end{corollary}

% =====================================================================
\section{Veredicto del programa}\label{sec:verdict}
% =====================================================================

\begin{enumerate}[leftmargin=1.6em]
\item \textbf{Densidad o posici\'on:} \textbf{posici\'on}. La cota
$\varepsilon_0\ge c/\lambda^2$ \emph{relevante para RH} es la del operador de Weil
$\Aweil$, y $\Aweil\ge0\Leftrightarrow$ RH. La cota del frame $\Sframe$ (densidad) es
incondicional pero \emph{no} es RH.
\item \textbf{Lo que el programa S\'I logr\'o} (incondicional, real):
\begin{itemize}[leftmargin=1.3em]
\item \PR\ Eslab\'on 1 (real--rootedness finita, H1+H2): $\xi$-truncadas con s\'olo
ceros reales, sin suponer RH.
\item \PR\ El motor CCM fiel (ceros a $10^{-50}$ con pocos primos).
\item \PR\ $\varepsilon_0^{\Sframe}\sim C_0/\lambda^2$ (densidad cr\'itica,
Landau--Beurling): un teorema de muestreo limpio.
\item \PR\ La anatom\'ia del margen: torre $(k+1)^2$, $\mathcal L$ rango--1,
direcci\'on--polo --- estructura del frame, no de RH.
\end{itemize}
\item \textbf{Lo que queda fuera de alcance no circular:} RH misma. El \'ultimo paso
($\Aweil\ge0$, o $\Aweil=\Sframe$) es el criterio de positividad de Weil, RH por
definici\'on. \emph{El programa es una reformulaci\'on limpia y una fuente de resultados
incondicionales de densidad, pero no un puente no circular a RH.}
\end{enumerate}

\noindent\textbf{Honestidad final.} Este es el desenlace esperable de cualquier
ataque a RH v\'ia positividad de Weil: el criterio es RH--equivalente, y separar
densidad de posici\'on muestra que la parte demostrable (densidad/frame) no toca la
parte RH (posici\'on). El valor del programa no est\'a en cerrar RH --- no la cierra ---
sino en (i)~los resultados incondicionales de los Eslabones 0,1 y del frame, y
(ii)~haber localizado, con demostraci\'on, que el muro es $\Aweil-\Sframe=$ posici\'on
$=$ RH, y no un operador $\mathcal L$ por descubrir ni una cola de Euler por acotar.
\end{document}
